\documentclass[12pt,a4paper]{article}
\usepackage[utf8]{inputenc}
\usepackage{amsmath,amssymb,amsthm}
\usepackage{geometry}
\usepackage{setspace}
\usepackage{hyperref}
\usepackage{natbib}
\geometry{margin=1in}
\onehalfspacing

\newtheorem{definition}{Definition}
\newtheorem{lemma}{Lemma}
\newtheorem{proposition}{Proposition}
\newtheorem{theorem}{Theorem}
\newtheorem{remark}{Remark}
\newtheorem{corollary}{Corollary}
\newtheorem{assumption}{Assumption}

\title{Cherry-Picking under Screening:\\
  Selective Disclosure and Participation in
  Pharmaceutical Approval}
\author{Xueyang}
\date{\today}

\begin{document}
\maketitle

% ============================================================
\section{Introduction}
% ============================================================

A pharmaceutical firm develops a drug and conducts a clinical
trial to persuade the FDA to approve it. The FDA does not know
the firm's type---the set of experiments available to the
firm---and designs an approval rule to screen different types.
This screening problem is studied by Min (2023), where the FDA
observes both the experiment chosen and its outcome.

I modify Min's model in two respects.

\emph{First (disclosure choice):} the trial outcome is the
firm's private information, and the firm chooses whether to
include it in its drug application. If included, the FDA
observes verified evidence. If omitted, the FDA observes only
an application without supporting trial data.

\emph{Second (participation):} drug development is costly.
The firm chooses whether to develop at all before running any
experiment. If the expected payoff from the approval mechanism
is too low, the firm does not develop---and no drug exists for
the FDA to evaluate.

\paragraph{Main results.}
\begin{enumerate}
  \item \textbf{Without the participation constraint}
    (Theorem~\ref{thm:no_IR}): the FDA's optimal committed
    mechanism always sets $A(\emptyset) = 0$ (reject
    applications without trial evidence). The model reduces
    exactly to Min (2023). Cherry-picking is
    payoff-irrelevant under commitment.

  \item \textbf{With the participation constraint}
    (Theorem~\ref{thm:with_IR}): when the small type's
    expected approval probability is too low to cover
    development costs, the FDA may optimally set
    $A(\emptyset) > 0$ to encourage participation. This
    $A(\emptyset) > 0$ arises \emph{endogenously} from the
    model---not from external assumptions.

  \item \textbf{Welfare cost of pooling}
    (Corollary~\ref{cor:welfare}): the pooling constraint
    ($A(\emptyset)$ is shared across all types) makes
    participation-inducing strictly more costly than in Min,
    where the FDA can set type-specific bad-result approval
    probabilities. The welfare difference measures the cost
    of cherry-picking.
\end{enumerate}


% ============================================================
\section{Model Setup}
% ============================================================

\subsection{Primitives}

Two players: a pharmaceutical firm (the agent, ``he'') and
the FDA (the principal, ``she'').

\paragraph{States and prior.}
$\omega \in \{G, B\}$, with
$\Pr(G) = p \in (0, 1/2)$.

\paragraph{Payoffs.}
The FDA receives 1 for a correct decision and 0 otherwise.
The firm cares only about approval:
$U_A(A, \cdot) = 1$, $U_A(D, \cdot) = 0$.

\paragraph{Development cost.}
Before running any experiment, the firm incurs a development
cost $c > 0$. If the firm does not develop, no drug exists
and the FDA makes no decision. The FDA's payoff from
non-participation equals $1 - p$ (equivalent to always
rejecting: correct on bad drugs, incorrect on good ones).

\begin{remark}[Role of $c$]
In Min (2023), the firm always participates
($c = 0$ implicitly). With $c > 0$, the firm's individual
rationality (IR) constraint can bind, creating a force for
$A(\emptyset) > 0$ that is absent from Min.
\end{remark}


% ------------------------------------------------------------
\subsection{Types and Feasible Sets}
% ------------------------------------------------------------

The firm has a privately known type
$\theta_i \in \{\theta_1, \theta_2\}$, with
$\Pr(\theta_1) = t$. Type determines the set of feasible
experiments $S_i$.

\begin{assumption}[Feasible sets]\label{as:feasible}
$S_2 \subset S_1$, and both satisfy:
\begin{enumerate}
  \item[(i)] Each $S_i$ is a closed, convex subset of
    $\{(x, y) \in [0,1]^2 : y \geq 1 - x\}$.
  \item[(ii)] The upper boundary
    $\bar{S}_i(x) := \max\{y : (x,y) \in S_i\}$ is
    continuous, concave, and strictly decreasing.
\end{enumerate}
\end{assumption}

Here $\theta_1$ is the ``big type'' (larger feasible set)
and $\theta_2$ is the ``small type.''


% ------------------------------------------------------------
\subsection{Experiments}
% ------------------------------------------------------------

An experiment $\pi = (x, y) \in S_i$ produces a binary outcome
$j \in \{g, b\}$:
\[
  \Pr(g \mid G, \pi) = y, \quad
  \Pr(g \mid B, \pi) = 1 - x.
\]
An experiment is \emph{convincing} if
$py \geq (1 - p)(1 - x)$.


% ============================================================
\subsection{Disclosure and Information Structure}
\label{sec:modification}
% ============================================================

After observing the trial outcome $j$, the firm submits a
drug application to the FDA. The firm chooses what evidence to
include:
\begin{itemize}
  \item \textbf{Disclose:} include the trial design and
    verified outcome $(\pi, j)$ in the application.
  \item \textbf{No evidence:} submit the application without
    trial data ($\emptyset$).
\end{itemize}

When the firm discloses, the FDA observes $(\pi, j)$, exactly
as in Min. When the firm submits without evidence, the FDA
sees only $\emptyset$: it does not learn which experiment was
conducted or what the outcome was.


% ============================================================
\subsection{Timeline}
\label{sec:timeline}
% ============================================================

\begin{enumerate}
  \item The FDA commits to approval probabilities:
    $A(\pi, j)$ for disclosed evidence, $A(\emptyset)$ for
    no-evidence applications.
  \item The firm learns its type $\theta_i$.
  \item The firm decides whether to develop (pay $c$) or
    not.
  \item If the firm develops: it conducts experiment
    $\pi \in S_i$, observes $j \in \{g, b\}$.
  \item The firm submits a drug application, choosing
    whether to include $(\pi, j)$ or submit with
    $\emptyset$.
  \item The FDA decides according to its committed rule.
\end{enumerate}

\begin{remark}[Comparison with Min]
Min's timeline has no step~3 (development decision) and no
step~5 (disclosure choice). Adding step~3 introduces the
participation constraint; adding step~5 introduces the
pooling constraint on $A(\emptyset)$.
\end{remark}


% ============================================================
\section{Analysis of the Disclosure Subgame}
\label{sec:disclosure}
% ============================================================

Conditional on the firm developing, the disclosure subgame
analysis is identical to the model without participation.

\begin{lemma}[Bad news]\label{lem:bad_news}
For any $\pi = (x,y)$ with $y \geq 1 - x$ and $p < 1/2$:
$\mu(G \mid \pi, b) < 1/2$.
\end{lemma}

\begin{proof}
$p(1-y) \leq p \cdot x < (1-p) \cdot x$.
\end{proof}

\begin{remark}[Silence dominates bad news informationally]
\label{rem:silence}
In any strategy profile, the FDA's posterior upon seeing
$\emptyset$ is weakly higher than upon seeing $b$:
\[
  \mu(G \mid \emptyset) \;\geq\; \mu(G \mid \pi, b),
\]
with strict inequality whenever silence can arise from a
hidden good result.

\medskip
\noindent\textit{Proof.}
Let $\sigma_g := \Pr(\text{hide}\mid g)$ and
$\sigma_b := \Pr(\text{hide}\mid b)$.
Write $H_g$ (resp.\ $H_b$) for the event that the firm
hides a good (resp.\ bad) result. Then
$\Pr(\emptyset) = \Pr(H_g) + \Pr(H_b)$, and by the law of
total probability:
\[
  \mu(G\mid\emptyset)
  \;=\;
  w\,\mu(G\mid g) \;+\; (1-w)\,\mu(G\mid b),
  \qquad
  w \;:=\; \frac{\Pr(H_g)}{\Pr(\emptyset)} \;\in\; [0,1].
\]
It remains to show $\mu(G\mid g) \geq \mu(G\mid b)$ for
any experiment $(x,y)$ with $y\geq 1-x$.
Cross-multiplying the inequality
$\frac{py}{py+(1-p)(1-x)} \geq \frac{p(1-y)}{p(1-y)+(1-p)x}$
reduces to
\[
  p(1-p)\,xy \;\geq\; p(1-p)(1-x)(1-y),
\]
i.e.\ $x+y \geq 1$, which holds by the feasibility
constraint $y\geq 1-x$.
Hence $\mu(G\mid\emptyset) \geq \mu(G\mid b)$ for any
$w\in[0,1]$, with strict inequality iff $w>0$, i.e.\
iff $\sigma_g>0$ (silence is genuinely ambiguous because
it can arise from hiding a good result).
\hfill$\square$

\medskip
The remark gives information-theoretic grounding for
$A(\emptyset)\geq A(b)$: silence is weakly better news
than an explicit bad result, so the FDA has no reason to
reward $b$ more than silence.
\end{remark}

\begin{lemma}[Equilibrium disclosure]\label{lem:disclosure}
The firm discloses $j$ iff $A(\pi, j) \geq A(\emptyset)$;
hides $j$ iff $A(\pi, j) < A(\emptyset)$.
\end{lemma}

\begin{proof}
The firm's payoff is $1$ if approved, $0$ otherwise.
Disclosing $j$ yields approval probability $A(\pi,j)$;
hiding $j$ yields $A(\emptyset)$, since the FDA sees only
$\emptyset$ and cannot distinguish which outcome was hidden.
The firm strictly prefers whichever action gives higher
approval probability, and is indifferent when
$A(\pi,j) = A(\emptyset)$.
\end{proof}

\begin{lemma}[WLOG: $A_b^i \leq A(\emptyset)$]
\label{lem:wlog_Ab}
In any optimal mechanism, $A(\pi_i, b) \leq A(\emptyset)$
for each type~$i$.
\end{lemma}

\begin{proof}
The FDA is the mechanism designer and strictly prefers low
$A(\pi_i, b)$: by Lemma~\ref{lem:bad_news},
$\mu(G \mid \pi_i, b) < 1/2$, so approving upon observing
$b$ is a mistake in expectation, and the FDA's ex-ante
payoff is strictly decreasing in $A(\pi_i, b)$.

It remains to show that the FDA can freely set
$A(\pi_i, b) \leq A(\emptyset)$ without affecting the
firm's payoff. By Lemma~\ref{lem:disclosure}, whenever
$A(\pi_i, b) < A(\emptyset)$ the firm hides $b$ and
receives $A(\emptyset)$. Hence the firm's effective payoff
from the $b$-outcome is $A(\emptyset)$ regardless of
$A(\pi_i, b) \leq A(\emptyset)$: the firm can always
guarantee $A(\emptyset)$ by hiding, so $A(\pi_i, b)$ below
this threshold is irrelevant to every IC and IR constraint.
The FDA therefore sets $A(\pi_i, b) \leq A(\emptyset)$,
and by Lemma~\ref{lem:disclosure} the firm hides $b$.
\end{proof}

\begin{lemma}[WLOG: $A_g^i \geq A(\emptyset)$]
\label{lem:wlog_Ag}
In any optimal mechanism, $A_g^i \geq A(\emptyset)$ for each
type~$i$.
\end{lemma}

\begin{proof}
If $A_g^i < A(\emptyset)$, the firm hides $g$ by
Lemma~\ref{lem:disclosure}, so both outcomes are hidden and
the experiment is payoff-irrelevant. Replacing $A_g^i$ with
$A(\emptyset)$ leaves all payoffs unchanged; then raising
$A_g^i$ above $A(\emptyset)$ for convincing experiments
weakly improves the FDA's payoff while satisfying
Lemma~\ref{lem:bad_news} for $g$ (since
$\mu(G \mid \pi, g) \geq 1/2$ for a convincing experiment).
\end{proof}

\begin{corollary}[Simplified disclosure structure]
\label{cor:disclosure}
In any optimal mechanism: the firm discloses $g$ and hides
$b$. The three lemmas above justify reducing the full
approval rule $\{A(\pi,g),\,A(\pi,b),\,A(\emptyset)\}$ to
the two-parameter object $\{A_g^i,\,A(\emptyset)\}$.
\end{corollary}

The simplified menu is:
\begin{equation}\label{eq:menu}
  \mathcal{M} = \left\{
    \left(\pi_1, A_g^1\right),\;
    \left(\pi_2, A_g^2\right);\;
    A(\emptyset)
  \right\}.
\end{equation}


% ============================================================
\section{The Mechanism Design Problem}
\label{sec:mechanism}
% ============================================================

\subsection{Payoffs}

Write $h_i := \Pr(g \mid \pi_i) = p y_i + (1-p)(1-x_i)$.

The firm's expected payoff from option $i$ (conditional on
developing):
\begin{equation}\label{eq:EU_A}
  EU^A_i = A_g^i \cdot h_i
    + A(\emptyset) \cdot (1 - h_i).
\end{equation}

The FDA's expected payoff from type $\theta_i$
(conditional on type~$i$ participating):
\begin{equation}\label{eq:EU_P_compact}
  EU^P_i = (1-p) + \Delta_g(\pi_i) \, A_g^i
           + \Delta_b(\pi_i) \, A(\emptyset),
\end{equation}
where:
\begin{align}
  \Delta_g(\pi_i)
    &:= p \, y_i - (1-p)(1-x_i),
    \label{eq:delta_g_def} \\
  \Delta_b(\pi_i)
    &:= p(1-y_i) - (1-p) x_i.
    \label{eq:delta_b_def}
\end{align}

Key identity: $\Delta_g + \Delta_b = 2p - 1 < 0$.


\subsection{Constraints}

\paragraph{Incentive compatibility.}
Type $\theta_1$ prefers option 1:
\begin{equation}\label{eq:IC1}
  \text{IC}_1: \quad
  A_g^1 h_1 + A(\emptyset)(1 - h_1)
  \geq
  A_g^2 h_2 + A(\emptyset)(1 - h_2).
\end{equation}
Type $\theta_2$ prefers option 2 (automatic if
$\pi_1 \notin S_2$).

\paragraph{Individual rationality (participation).}
Type $\theta_i$ develops iff:
\begin{equation}\label{eq:IR}
  \text{IR}_i: \quad
  A_g^i \cdot h_i + A(\emptyset) \cdot (1 - h_i) \geq c.
\end{equation}


\subsection{FDA's Problem}

Define the indicator $\mathbb{1}_i = 1$ if type~$i$
participates, $0$ otherwise. The FDA's total expected payoff:
\begin{equation}\label{eq:fda_total}
  V = \sum_{i=1}^{2} \Pr(\theta_i) \left[
    \mathbb{1}_i \cdot EU^P_i
    + (1 - \mathbb{1}_i) \cdot (1-p)
  \right].
\end{equation}

Since $EU^P_i \geq 1-p$ for convincing experiments with
$A_g^i > 0$ (because $\Delta_g A_g^i > 0$ and
$\Delta_b A(\emptyset) < 0$, but net is positive for
reasonable mechanisms), the FDA generally prefers
participation.

The FDA maximizes $V$ over
$(\pi_1, \pi_2, A_g^1, A_g^2, A(\emptyset))$ subject to
IC, IR, and feasibility.


% ============================================================
\section{Main Results}
\label{sec:main}
% ============================================================

\subsection{Without Participation Constraint}

\begin{theorem}[Irrelevance of cherry-picking under
  commitment]\label{thm:no_IR}
If $\mathrm{IR}_i$ does not bind for either type (e.g.,
$c = 0$), then the optimal mechanism has
$A(\emptyset) = 0$. The FDA's optimal payoff equals Min
(2023)'s optimal payoff.
\end{theorem}

\begin{proof}
Fix any experiments $\pi_1, \pi_2$. The FDA's payoff is:
\[
  V = (1-p)
    + t \, \Delta_g(\pi_1) A_g^1
    + (1-t) \Delta_g(\pi_2) A_g^2
    + \bigl[t \, \Delta_b(\pi_1)
      + (1-t) \Delta_b(\pi_2)\bigr] A(\emptyset).
\]
We show $dV/dA(\emptyset) < 0$ in all cases.

\medskip
\noindent\textbf{Case 1: No IC binds.}
$A_g^i$ do not adjust with $A(\emptyset)$:
$dV/dA(\emptyset) = t \Delta_b(\pi_1)
+ (1-t) \Delta_b(\pi_2) < 0$.

\medskip
\noindent\textbf{Case 2: $\mathrm{IC}_1$ binds,
$A_g^1 = 1$.}
Binding $\mathrm{IC}_1$:
$A_g^2 = [h_1 + A(\emptyset)(h_2 - h_1)]/h_2$.
Then $dA_g^2/dA(\emptyset) = (h_2 - h_1)/h_2$.

\begin{align*}
  \frac{dV}{dA(\emptyset)}
  &= (1-t)\left[
       \Delta_g(\pi_2) \cdot \frac{h_2 - h_1}{h_2}
       + \Delta_b(\pi_2)
     \right]
     + t \, \Delta_b(\pi_1) \\
  &= (1-t)\left[
       \underbrace{(\Delta_g(\pi_2) + \Delta_b(\pi_2))}
         _{= \, 2p-1 \, < \, 0}
       - \Delta_g(\pi_2) \cdot \frac{h_1}{h_2}
     \right]
     + t \, \Delta_b(\pi_1) < 0.
\end{align*}

\medskip
\noindent\textbf{Case 3: $\mathrm{IC}_2$ binds,
$A_g^2 = 1$.}
By symmetric reasoning:
$dV/dA(\emptyset) = t[(2p-1) - \Delta_g(\pi_1) h_2/h_1]
+ (1-t) \Delta_b(\pi_2) < 0$.

\medskip
Since $dV/dA(\emptyset) < 0$ for all experiment pairs,
$A(\emptyset) = 0$ is globally optimal. At
$A(\emptyset) = 0$: $EU^A = A_g^i h_i$, identical to Min
with $A_b^i = 0$.
\end{proof}

\begin{remark}[Economic intuition]
The key identity is $\Delta_g + \Delta_b = 2p - 1 < 0$.
Even when raising $A(\emptyset)$ relaxes an IC constraint
(the screening benefit), this benefit equals at most
$\Delta_g$. The direct cost is $|\Delta_b|$. Since
$\Delta_g < |\Delta_b|$ (because their sum is $2p-1 < 0$),
the cost always exceeds the benefit.

Put differently: $A(\emptyset)$ is a \textbf{pooled}
instrument affecting all types. Min's $A_b^i$ is
type-specific. The pooling makes $A(\emptyset)$ too
blunt---it always does more harm than good as a screening
tool.
\end{remark}


% ============================================================
\subsection{With Participation Constraint}
\label{sec:with_IR}
% ============================================================

We now study the case where IR can bind. Focus on the
leading case: $\pi_1 \notin S_2$ (so $\mathrm{IC}_2$ is
automatic and the small type cannot mimic the big type's
experiment).

\begin{lemma}[When does IR bind?]\label{lem:IR_binding}
Suppose $\pi_1 \notin S_2$ and $\mathrm{IC}_1$ does not
bind. Then the optimal mechanism sets $A_g^1 = A_g^2 = 1$,
and type~$i$'s payoff is:
\[
  EU^A_i = h_i + A(\emptyset)(1 - h_i).
\]
At $A(\emptyset) = 0$: $EU^A_i = h_i$. Since
$h_2 \leq h_1$ (because $S_2 \subset S_1$), type~2 is
the marginal participant.

If $\mathrm{IC}_1$ binds (binding at $A(\emptyset) = 0$),
then type~2's payoff equals type~1's payoff $= h_1 \geq c$
whenever type~1 participates. In this case, $\mathrm{IR}_2$
is automatically satisfied.
\end{lemma}

\begin{proof}
When $\mathrm{IC}_1$ does not bind, $A_g^i$ are
unconstrained and the FDA sets $A_g^i = 1$ (optimal for
convincing experiments). Type~$i$'s payoff is
$h_i + A(\emptyset)(1-h_i)$, which is $h_i$ at
$A(\emptyset) = 0$.

When $\mathrm{IC}_1$ binds with $A_g^1 = 1$:
$h_1 + A(\emptyset)(1-h_1)
= A_g^2 h_2 + A(\emptyset)(1-h_2)$.
Type~2's equilibrium payoff from option~2 equals
type~1's payoff from option~1 $= h_1 + A(\emptyset)(1-h_1)$.
At $A(\emptyset) = 0$: type~2's payoff $= h_1 \geq c$.
\end{proof}

\begin{remark}
Lemma~\ref{lem:IR_binding} says the interesting case is:
$\mathrm{IC}_1$ slack and $h_1 \geq c > h_2$. The big type
can profitably develop; the small type cannot, unless the
FDA provides a participation subsidy through
$A(\emptyset) > 0$.
\end{remark}

\begin{theorem}[Endogenous $A(\emptyset) > 0$]
\label{thm:with_IR}
Suppose $\pi_1 \notin S_2$, $\mathrm{IC}_1$ is slack,
$A_g^1 = A_g^2 = 1$, and $h_1 \geq c > h_2$.

Define $\alpha^* := (c - h_2)/(1 - h_2)$, the minimum
$A(\emptyset)$ that induces type~2's participation. Then
$A(\emptyset) = \alpha^*$ is optimal if and only if:
\begin{equation}\label{eq:threshold_condition}
  (1-t) \, \Delta_g(\pi_2)
  > \bigl|t \, \Delta_b(\pi_1)
    + (1-t) \, \Delta_b(\pi_2)\bigr|
    \cdot \alpha^*.
\end{equation}
\end{theorem}

\begin{proof}
\emph{Step~1: Structure of $V(A(\emptyset))$.}
At $A(\emptyset) = \alpha$:
\begin{itemize}
  \item If $\alpha < \alpha^*$: type~2 does not participate.
    $V(\alpha) = t[(1-p) + \Delta_g(\pi_1)
    + \Delta_b(\pi_1) \alpha] + (1-t)(1-p)$.
    This is linear and decreasing in $\alpha$ (slope
    $= t \, \Delta_b(\pi_1) < 0$).
  \item If $\alpha \geq \alpha^*$: both types participate.
    $V(\alpha) = (1-p)
    + t[\Delta_g(\pi_1) + \Delta_b(\pi_1) \alpha]
    + (1-t)[\Delta_g(\pi_2) + \Delta_b(\pi_2) \alpha]$.
    This is linear and decreasing in $\alpha$ (slope
    $= t \Delta_b(\pi_1) + (1-t) \Delta_b(\pi_2) < 0$
    by Theorem~\ref{thm:no_IR}).
\end{itemize}
So $V$ is piecewise linear: decreasing on
$[0, \alpha^*)$, then an upward jump at $\alpha^*$
(from type~2 joining), then decreasing on
$[\alpha^*, 1]$.

\medskip
\emph{Step~2: The jump at $\alpha^*$.}
\begin{align*}
  V(\alpha^*) - V(\alpha^{*-})
  &= (1-t) \left[
       \Delta_g(\pi_2) + \Delta_b(\pi_2) \alpha^*
     \right] \\
  &= (1-t) \left[
       \Delta_g(\pi_2)(1 - \alpha^*)
       + (2p-1) \alpha^*
     \right].
\end{align*}
For the jump to be positive, we need
$\Delta_g(\pi_2)(1-\alpha^*) > (1-2p)\alpha^*$, i.e.,
$\Delta_g(\pi_2) + \Delta_b(\pi_2) \alpha^* > 0$. This
is the participation value condition: the FDA's gain from
type~2's good-result approvals exceeds the cost from
type~2's silence-channel approvals.

\medskip
\emph{Step~3: Comparison with $\alpha = 0$.}
Since $V$ is decreasing on $[0, \alpha^*)$, the best
below the jump is $\alpha = 0$. Since $V$ is decreasing
on $[\alpha^*, 1]$, the best above the jump is
$\alpha = \alpha^*$. So the global optimum is either
$\alpha = 0$ or $\alpha = \alpha^*$.
\begin{align}
  V(\alpha^*) - V(0)
  &= (1-t) \, \Delta_g(\pi_2)
     + \bigl[t \, \Delta_b(\pi_1)
       + (1-t) \, \Delta_b(\pi_2)\bigr] \alpha^*.
     \label{eq:net_gain}
\end{align}
This is positive iff
\eqref{eq:threshold_condition} holds.
\end{proof}

\begin{corollary}[Development cost threshold]
\label{cor:threshold}
Define:
\begin{equation}\label{eq:cstar}
  c^* := h_2
  + \frac{(1-t) \, \Delta_g(\pi_2) \cdot (1 - h_2)}
         {\bigl|t \, \Delta_b(\pi_1)
           + (1-t) \, \Delta_b(\pi_2)\bigr|}.
\end{equation}
Then:
\begin{enumerate}
  \item[(i)] If $c \leq h_2$: both types participate at
    $A(\emptyset) = 0$. By Theorem~\ref{thm:no_IR},
    $A(\emptyset) = 0$ optimal. No cherry-picking.
  \item[(ii)] If $c \in (h_2, c^*)$: type~2 does not
    participate at $A(\emptyset) = 0$. The FDA optimally
    sets $A(\emptyset) = \alpha^* = (c - h_2)/(1-h_2) > 0$
    to induce participation. \textbf{Cherry-picking
    is endogenous.}
  \item[(iii)] If $c \geq c^*$: the cost of
    participation-inducing $A(\emptyset)$ exceeds the
    benefit. Type~2 does not participate; $A(\emptyset) = 0$.
\end{enumerate}
\end{corollary}

\begin{proof}
In \eqref{eq:net_gain}, substitute
$\alpha^* = (c - h_2)/(1-h_2)$ and solve
$V(\alpha^*) - V(0) > 0$ for $c$. This gives $c < c^*$.
\end{proof}

\begin{remark}[When is $A(\emptyset) > 0$ most likely
  optimal?]
Condition \eqref{eq:threshold_condition} is most easily
satisfied when:
\begin{itemize}
  \item $(1-t)$ is large: the small type is common, so
    losing its participation is costly.
  \item $\Delta_g(\pi_2)$ is large: the small type's
    experiment is highly convincing, so its participation
    is valuable to the FDA.
  \item $\alpha^*$ is small: the participation constraint
    barely binds ($c$ close to $h_2$), so little
    $A(\emptyset)$ subsidy is needed.
\end{itemize}
\end{remark}


% ============================================================
\subsection{Welfare Cost of Pooling}
\label{sec:welfare}
% ============================================================

\begin{corollary}[Welfare cost of cherry-picking]
\label{cor:welfare}
Consider the same model primitives
$(p, t, S_1, S_2, c)$ under Min's full-transparency regime
(FDA observes $(\pi, j)$ always) and under my voluntary
disclosure regime. Both have the participation constraint
$EU^A_i \geq c$.

When $c \in (h_2, c^*)$:
\begin{enumerate}
  \item[(i)] \textbf{Min:} the FDA sets $A_b^2 > 0$
    specifically for type~2, leaving $A_b^1 = 0$. Cost:
    only type~2's bad-result channel is opened.
  \item[(ii)] \textbf{My model:} the FDA sets
    $A(\emptyset) = \alpha^* > 0$ for all types. Cost:
    both types' silence channels are opened.
\end{enumerate}
The welfare difference is:
\begin{equation}\label{eq:welfare_cost}
  V_{\text{Min}} - V_{\text{mine}}
  = t \cdot |\Delta_b(\pi_1)| \cdot \alpha^* > 0.
\end{equation}
This is the cost of the \textbf{pooling constraint}:
the FDA must subsidize type~1 (who doesn't need it)
in order to subsidize type~2 (who does).
\end{corollary}

\begin{proof}
In Min, $\mathrm{IR}_2$ requires
$A_g^2 h_2 + A_b^2 (1-h_2) \geq c$. With $A_g^2 = 1$:
$A_b^2 = (c - h_2)/(1-h_2) = \alpha^*$. Set $A_b^1 = 0$.

Min's payoff:
$V_{\text{Min}} = t[(1-p) + \Delta_g(\pi_1)]
+ (1-t)[(1-p) + \Delta_g(\pi_2) + \Delta_b(\pi_2) \alpha^*]$.

My payoff:
$V_{\text{mine}} = t[(1-p) + \Delta_g(\pi_1)
+ \Delta_b(\pi_1) \alpha^*]
+ (1-t)[(1-p) + \Delta_g(\pi_2) + \Delta_b(\pi_2) \alpha^*]$.

Difference: $V_{\text{Min}} - V_{\text{mine}}
= -t \, \Delta_b(\pi_1) \alpha^*
= t \, |\Delta_b(\pi_1)| \, \alpha^* > 0$.
\end{proof}

\begin{remark}[Interpretation]
In Min, $A_b^2 > 0$ is a targeted subsidy: only type~2
gets bad-result approval. In my model, $A(\emptyset)$ is a
blanket subsidy: all types get silence approval. The excess
cost $t |\Delta_b(\pi_1)| \alpha^*$ is the price of
type~1's unnecessary cherry-picking channel.

This is the \textbf{welfare cost of cherry-picking}: not
from the act of hiding per se, but from the FDA's inability
to target bad-result tolerance to specific types.
\end{remark}


% ============================================================
\section{Discussion}
\label{sec:discussion}
% ============================================================

\subsection{Summary of Results}

\begin{center}
\renewcommand{\arraystretch}{1.3}
\begin{tabular}{l|c|c}
& \textbf{$c \leq h_2$} & \textbf{$c \in (h_2, c^*)$} \\
\hline
$A(\emptyset)$ & $0$ & $\alpha^* > 0$ \\
Cherry-picking & irrelevant & endogenous \\
Equiv.\ to Min? & yes & no \\
Welfare cost & $0$ & $t|\Delta_b(\pi_1)|\alpha^*$
\end{tabular}
\end{center}


\subsection{Limitations and Honest Assessment}

\paragraph{What is rigorous:}
\begin{itemize}
  \item Lemmas 1--3 (bad news, disclosure, WLOG $A_g$):
    straightforward, fully proved.
  \item Theorem~\ref{thm:no_IR} ($A(\emptyset) = 0$
    without IR): fully proved for all IC cases.
  \item Theorem~\ref{thm:with_IR} (endogenous
    $A(\emptyset) > 0$): fully proved for the case
    $\pi_1 \notin S_2$, $\mathrm{IC}_1$ slack.
  \item Corollary~\ref{cor:welfare} (welfare cost):
    clean comparison under the same assumptions.
\end{itemize}

\paragraph{What requires further work:}
\begin{itemize}
  \item When $\pi_1 \in S_2$ (both IC constraints active):
    the interaction between IC and IR is more complex.
    Closed-form characterization is not available; this is
    where numerical methods are needed.
  \item The threshold $c^*$ depends on fixed experiments
    $\pi_1, \pi_2$. Joint optimization over experiments
    and approval probabilities with IR is a
    higher-dimensional problem without clean closed form.
  \item Min's Propositions~4--5 (sufficient conditions for
    favorite experiments, ex post optimality) have not been
    extended to the IR case. This extension is non-trivial.
\end{itemize}


\subsection{Interpreting $A(\emptyset) > 0$}

The model's prediction that the FDA optimally commits to
$A(\emptyset) > 0$ maps to several real-world regulatory
institutions:
\begin{itemize}
  \item \textbf{Accelerated/conditional approval:}
    The FDA approves drugs based on preliminary or surrogate
    evidence, requiring post-market confirmatory studies.
    The firm submits with incomplete data ($\emptyset$) and
    receives approval with positive probability---exactly
    the participation subsidy in the model.

  \item \textbf{Emergency use authorization:}
    EUAs allow approval before full trial completion, a
    direct analog of $A(\emptyset) > 0$ when development
    incentives are especially pressing.

  \item \textbf{Silence as preliminary outcome:}
    From the FDA's perspective, $\emptyset$ is
    informationally ambiguous---it could reflect a hidden
    bad result \emph{or} an ongoing trial with no
    conclusion yet. By Remark~\ref{rem:silence},
    $\mu(G|\emptyset) \geq \mu(G|b)$: silence is weakly
    better news than confirmed bad results. The FDA's
    willingness to set $A(\emptyset) > 0$ reflects this
    ambiguity.
\end{itemize}
The key insight: $A(\emptyset) > 0$ is not regulatory
negligence but a rational design choice---the FDA accepts
higher risk of approving bad drugs to ensure good drugs get
developed. Corollary~\ref{cor:welfare} quantifies the cost.


\subsection{Policy Implications}

\paragraph{Pre-registration as restoring targeting.}
Mandatory trial pre-registration makes the experiment
$\pi$ observable even when the firm omits results. This
breaks the pooling constraint:
$A(\emptyset) \to A(\pi, \emptyset)$ becomes
experiment-specific, like Min's $A_b^i$. The FDA can
target subsidies, eliminating the welfare cost of pooling.

\paragraph{Reducing development costs.}
Government subsidies or tax credits that lower $c$
can eliminate the need for $A(\emptyset) > 0$ entirely
(moving from regime (ii) to (i) in the table above).
This is more efficient than tolerating cherry-picking.


% ============================================================
\section*{References}
% ============================================================

\begin{itemize}
  \item Min, D.\ (2023). Screening for experiments.
    \emph{Games and Economic Behavior}, 142, 898--920.
\end{itemize}

\end{document}
